% Einheit 4 – Sachaufgaben – Nr. 45–49
\setcounter{aufgabe}{44}
\einheitenkopf[e4]{Einheit 4 von 4 · Sachaufgaben}
\zweigzeile{Hier lernst du, zu Zahlenrätseln und Rechtecken eine quadratische Gleichung aufzustellen, sie zu lösen und die passende Lösung auszuwählen · neu in diesem Jahr (an der Oberschule je nach Buch Kl. 9 oder 10) · keine P10-Aufgabe · baut auf: Einheit 1 bis 3, Gleichungen aufstellen}

\begin{aufgabe}{Ich erkenne, welche Lösung zur Frage passt.}
Rechne nicht. Eine Gleichung zur Frage hat die Lösungen $x_1$ und $x_2$. Kreuze an, welche Lösung als Antwort passt.
\begin{teile}
\teil Seitenlänge eines quadratischen Beets in m: $x_1 = 6$, $x_2 = -6$ \hfill \kreuz{$x_1$} \quad \kreuz{$x_2$} \quad \kreuz{beide}
\teil Breite eines Rechtecks in cm: $x_1 = -12$, $x_2 = 4$ \hfill \kreuz{$x_1$} \quad \kreuz{$x_2$} \quad \kreuz{beide}
\teil Anzahl der Stuhlreihen in einem Saal: $x_1 = 9$, $x_2 = -14$ \hfill \kreuz{$x_1$} \quad \kreuz{$x_2$} \quad \kreuz{beide}
\teil Zahlen, deren Quadrat 144 ist: $x_1 = 12$, $x_2 = -12$ \hfill \kreuz{$x_1$} \quad \kreuz{$x_2$} \quad \kreuz{beide}
\teil Alter eines Kindes in Jahren: $x_1 = -3$, $x_2 = 11$ \hfill \kreuz{$x_1$} \quad \kreuz{$x_2$} \quad \kreuz{beide}
\end{teile}
\end{aufgabe}

\verfahren{Zahlenrätsel}

\begin{aufgabe}{Ich kann ein Zahlenrätsel mit einer quadratischen Gleichung lösen.}
Löse das Rätsel. Gib alle passenden Zahlen an.
\begin{teile}
\teil Das Quadrat einer Zahl ist 100. \quad Gleichung: $x^2 = 100$ \hfill Zahlen: \leerfeld
\teil Das Quadrat einer Zahl ist 121. \quad Gleichung: $x^2 = 121$ \hfill Zahlen: \leerfeld
\teil Das Doppelte des Quadrats einer Zahl ist 72. \quad Gleichung: $2x^2 = 72$ \hfill Zahlen: \leerfeld
\teil Vermindert man das Quadrat einer Zahl um 10, erhält man 71. \\ Gleichung: $x^2 - 10 = 71$ \hfill Zahlen: \leerfeld
\anweisung{Stelle die Gleichung selbst auf.}
\teil Addiert man 20 zum Quadrat einer Zahl, erhält man 84. \\ Gleichung: \leerfeld \hfill Zahlen: \leerfeld
\teil Das Produkt einer Zahl und ihres Nachfolgers ist 56. \\ Gleichung: \leerfeld \hfill Zahlenpaare: \leerfeld
\end{teile}
\end{aufgabe}

\verfahren{Rechteckaufgaben}

\begin{aufgabe}{Ich kann eine Rechteckaufgabe mit einer quadratischen Gleichung lösen.}
Ein rechteckiges Beet ist 3\,m länger als breit und hat einen Flächeninhalt von $40\,\text{m}^2$.

\rechteck[punkte=,seiten={x + 3,x,,}]{4.8}{3}

\begin{teile}
\teil Stelle eine Gleichung für den Flächeninhalt auf. \hfill \leerfeld
\teil Löse die Gleichung. \hfill \feld{x_1} \quad \feld{x_2}
\teil Begründe, welche Lösung zur Aufgabe passt.
\teil Gib Länge und Breite des Beets in einem Antwortsatz mit Einheit an.
\teil Prüfe am Beet, ob der Flächeninhalt stimmt.
\end{teile}
Um ein anderes rechteckiges Beet, das 4\,m länger als breit ist, wird ringsum ein 1\,m breiter Weg angelegt. Beet und Weg bedecken zusammen $96\,\text{m}^2$.

\rechteckrand{2}{3.6}{0.6}{x}{x + 4}{Weg, 1\,m breit}

\begin{teile}
\teil Berechne die Länge und die Breite des Beets. \hfill \feld[m]{\text{Breite}} \quad \feld[m]{\text{Länge}}
\end{teile}
\end{aufgabe}

\begin{aufgabe}{Ich finde den Fehler bei einer Sachaufgabe.}
Finde den Fehler, benenne ihn und korrigiere.
\begin{teile}
\teil Ein Rechteck ist 2\,cm länger als breit und hat einen Flächeninhalt von $48\,\text{cm}^2$. Sven rechnet:
\rechnung{x\cdot(x + 2) &= 48 \\ x^2 + 2x - 48 &= 0 \\ x_1 = 6, \quad x_2 &= -8}
Svens Antwort: Das Rechteck ist $-8$\,cm breit und $-6$\,cm lang.
\teil Ein Rechteck ist 5\,m länger als breit und hat einen Flächeninhalt von $24\,\text{m}^2$. Lisa rechnet:
\rechnung{2x + 2\cdot(x + 5) &= 24 \\ 4x + 10 &= 24 \\ x &= 3{,}5}
\end{teile}
\end{aufgabe}

\begin{aufgabe}{Ich kann begründen, wann eine Sachaufgabe eine oder zwei Antworten hat.}
Begründe.
\begin{teile}
\teil Warum passt bei einer Rechteckaufgabe meist nur eine Lösung der Gleichung?
\teil Warum hat das Rätsel „Das Quadrat einer Zahl ist 25. Wie heißt die Zahl?“ zwei Antworten?
\end{teile}
\end{aufgabe}
